\section{Geometric and horizontal preparations}
\label{sec:preliminaries}

We use two established results: the pointwise Gaussian marginal theorem
of Eldan and Klartag, and the Rogers--Schmidt mean-hole estimate. From
the first we derive a finite list of bodies contained in nearby horizontal
sections; from the second we obtain one horizontal lattice that works for
the whole list.  Neither result requires symmetry. This section develops the finite-section
and common-lattice preparation of Li and Liu~\cite[Proposition~4.1
and Lemma~4.2]{LL}, with the precise bounds used in our construction.

\subsection{A finite list of horizontal sections}

Theorem~1 of Eldan and Klartag~\cite{EK} has the following consequence.
There are absolute constants $\alpha,\beta,\gamma,C_0>0$ such that, for
all sufficiently large $n$, every isotropic random vector $X\in\R^n$
with a log-concave density has the following property: for each integer
$1\le D\le n^\alpha$, there is a $D$-dimensional subspace $E$ whose
orthogonal marginal has density $f_E$ satisfying
\begin{equation}
 \left|\frac{f_E(y)}{\gamma_D(y)}-1\right|
 \le C_0 n^{-\gamma}
 \qquad (y\in E,\ \|y\|\le n^\beta).
 \label{eq:gaussian-marginal-input}
\end{equation}
Here isotropic means $\E X=0$ and $\E XX^{\mathsf T}=I_n$; orthonormal
coordinates identify $E$ with $\R^D$.  For a uniform law on a convex
body, we use the canonical marginal density given by the section integral.
We only use the
lower bound in Equation~\eqref{eq:gaussian-marginal-input}.

\begin{lemma}[Section labels]
\label{lem:sections}
There are absolute constants $c,C>0$ with the following property.
Fix $M\ge1$.  For all sufficiently large $n$, with threshold depending
only on $M$, suppose that $D$ is a positive integer and
\[
 1\le D,R\le M(1+\log\log n)^M,\qquad m=n-D.
\]
Every convex body in $\R^n$ has an invertible affine image
$K\subset\R^m\times\R^D$ and a list of at most $(Cm^2)^D$ convex
bodies $J_a\subset\R^m$ such that, writing $K_y=\{x:(x,y)\in K\}$,
for every $\|y\|\le R$ some label $a$ satisfies
\begin{equation}
 J_a\subset K_y,
 \qquad
 c\,\vol_n(K)\gamma_D(y)
 \le \vol_m(J_a)
 \le C\,\vol_n(K)\gamma_D(y).
 \label{eq:section-labels}
\end{equation}
\end{lemma}

\begin{proof}
An affine transformation makes the uniform law on the given body
isotropic: its covariance matrix is positive definite because the body
has nonempty interior.  The uniform density is log-concave.  Put
$S=2DR+2$.  Uniformly over the displayed range of $D,R$, for sufficiently
large $n$ we have
\[
 D\le n^\alpha,\qquad S<n^\beta,
 \qquad C_0n^{-\gamma}\le\tfrac12.
\]
Choose $E$ from Equation~\eqref{eq:gaussian-marginal-input} and use
orthonormal coordinates $E^\perp\times E$.  Write $K$ for the resulting
body.  By Fubini's theorem, $\vol_m(K_y)/\vol_n(K)$ is a version of
the marginal density.  The lower bound holds at every required height,
independently of density versions: compactness implies that, as $y'\to y$,
the fiber $K_{y'}$ is eventually contained in $K_y+\eps B_m$ for each
$\eps>0$, where $B_m$ is the unit ball.  Continuity of volume under
decreasing compact neighborhoods shows that $y\mapsto\vol_m(K_y)$ is
upper semicontinuous.  Its almost-everywhere Gaussian lower bound therefore
extends to every $\|y\|\le S<n^\beta$ by approaching from the full-measure
set.  Consequently
\begin{equation}
 \vol_m(K_y)\ge\tfrac12\vol_n(K)\gamma_D(y)
 \qquad (\|y\|\le S).
 \label{eq:section-lower-bound}
\end{equation}

We first arrange that $0\in K_y$ whenever $\|y\|\le2R$.
Take a centered regular simplex in $\R^D$ of inradius $2R$; its vertices
$v_0,\ldots,v_D$ have norm $2DR<S$.  Each section at a vertex is nonempty
by Equation~\eqref{eq:section-lower-bound}, so choose $x_i\in K_{v_i}$.
There is an affine map $T:\R^D\to\R^m$ with $T(v_i)=x_i$ for all $i$.
Replace the body by its image under $(x,y)\mapsto(x-T(y),y)$.
This map has determinant one and translates each horizontal section;
thus it preserves both its volume and Equation~\eqref{eq:section-lower-bound}.
The new body contains $(0,v_i)$ for every vertex.  By convexity it contains
$(0,y)$ for every $y$ in the simplex, which includes the ball of radius
$2R$.

Let $\{y_a\}$ be a maximal $R/m^2$-separated set in the closed ball of
radius $R$.  It is an $R/m^2$-net.  Comparing the disjoint balls of radius
$R/(2m^2)$ about its points with the ball of radius $R+R/(2m^2)$ gives
\begin{equation}
 \#\{y_a\}\le(1+2m^2)^D\le(3m^2)^D.
 \label{eq:section-net-size}
\end{equation}
If $\|y-y_a\|\le R/m^2$, set $z=y_a+m^2(y-y_a)$.  Then
$\|z\|\le2R$ and
$y=(1-m^{-2})y_a+m^{-2}z$.  Since $(0,z)\in K$, convexity yields
\begin{equation}
 (1-m^{-2})K_{y_a}\subset K_y.
 \label{eq:section-net-inclusion}
\end{equation}

For $m\ge2$, Bernoulli's inequality gives
$(1-m^{-2})^m\ge1-m^{-1}\ge1/2$.
By Equation~\eqref{eq:section-lower-bound}, the body
$(1-m^{-2})K_{y_a}$ therefore has volume at least
$\tfrac14\vol_n(K)\gamma_D(y_a)$.
It contains zero, so a further homothety about zero gives a convex body
$J_a$ contained in it with exactly
\begin{equation}
 \vol_m(J_a)=\tfrac18\vol_n(K)\gamma_D(y_a).
 \label{eq:label-prescribed-volume}
\end{equation}
For a chosen nearby $y_a$,
\[
 \left|\log\frac{\gamma_D(y_a)}{\gamma_D(y)}\right|
 \le \frac{\|y_a-y\|(\|y_a\|+\|y\|)}2
 \le R^2/m^2=o(1).
\]
Equations~\eqref{eq:section-net-inclusion} and
\eqref{eq:label-prescribed-volume} prove Equation~\eqref{eq:section-labels},
and Equation~\eqref{eq:section-net-size} bounds the list size.
All thresholds used above depend only on $M$, not on the body.
\end{proof}

\subsection{One lattice for all labels}

For a lattice $\Lambda\subset\R^m$, write
\[
 h_{\Lambda}(J)
 =1-|\pi_\Lambda(J)|,
\]
where the measure on $\R^m/\Lambda$ is normalized to be a probability
measure.  Let $\mu_m$ be the invariant probability measure on the space
of unimodular lattices in $\R^m$, and put
\[
 \eta_m=\frac m4\log\frac{27}{16}-3\log m.
\]
The estimate of Rogers and Schmidt~\cite{Rog58,Sch58}, in the form of
Ordentlich--Regev--Weiss~\cite[Theorem~2.3]{ORW}, states that there is an
absolute constant $C_{\mathrm R}>0$ such that every Borel set
$J\subset\R^m$ of volume $u\le\eta_m$ satisfies
\begin{equation}
 \left|\int h_{\Lambda}(J)\,d\mu_m(\Lambda)-e^{-u}\right|
 \le C_{\mathrm R}e^{-\eta_m}.
 \label{eq:mean-hole-input}
\end{equation}

\begin{lemma}[Common horizontal lattice]
\label{lem:horizontal}
Let $J_1,\ldots,J_N$ be convex bodies in $\R^m$, let $D_h>0$, and set
$u_a=\vol_m(J_a)/D_h$.  If
\begin{equation}
 \max_a u_a\le\eta_m,
 \qquad
 (1+C_{\mathrm R})\sum_{a=1}^N e^{-u_a/2}<1,
 \label{eq:horizontal-criterion}
\end{equation}
there is a lattice $\Lambda_h$ of determinant $D_h$ satisfying
\begin{equation}
 h_{\Lambda_h}(J_a)\le e^{-u_a/2}
 \qquad (1\le a\le N).
 \label{eq:horizontal-holes}
\end{equation}
In particular, this conclusion holds for all sufficiently large $m$ if
\[
 \min_a u_a\ge\sqrt m,\qquad
 \max_a u_a=o(m),\qquad
 \log N=o(\sqrt m).
\]
The dimension threshold in this last formulation depends only on the
three numerical bounds, not on the shapes of the bodies.
\end{lemma}

\begin{proof}
Replace every $J_a$ by $J'_a=D_h^{-1/m}J_a$, so that
$\vol_m(J'_a)=u_a$.  Equation~\eqref{eq:mean-hole-input} gives
\[
 \E_{\Lambda\sim\mu_m}h_{\Lambda}(J'_a)
 \le e^{-u_a}+C_{\mathrm R}e^{-\eta_m}
 \le(1+C_{\mathrm R})e^{-u_a}.
\]
Markov's inequality bounds the probability that
$h_{\Lambda}(J'_a)>e^{-u_a/2}$ by
$(1+C_{\mathrm R})e^{-u_a/2}$.  The sum of these probabilities is less
than one by Equation~\eqref{eq:horizontal-criterion}.  Thus a single
unimodular lattice $\Lambda$ satisfies all the desired inequalities.
Set $\Lambda_h=D_h^{1/m}\Lambda$; scaling identifies the normalized
tori, preserves the hole proportions, and gives determinant $D_h$.

Finally, $\eta_m/m\to\tfrac14\log(27/16)>0$, so
$\max_a u_a=o(m)$ ensures the first condition in
Equation~\eqref{eq:horizontal-criterion}.  The other two numerical
bounds give
\[
 (1+C_{\mathrm R})\sum_{a=1}^N e^{-u_a/2}
 \le(1+C_{\mathrm R})\exp\{\log N-\tfrac12\sqrt m\}=o(1),
\]
which ensures the second.
\end{proof}
